\documentclass[a4paper,12pt,twoside,final]{scrbook}
\usepackage[utf8]{inputenc}
\usepackage[spanish,es-nodecimaldot]{babel} % International, rec. by RAE: http://www.tex-tipografia.com/marca_decimal.html

\usepackage{booktabs}
\usepackage[hidelinks]{hyperref}

% COLOR
\usepackage[usenames,dvipsnames,svgnames,table]{xcolor}
% Colores para estilo Proyecto Docente (tonos naranjas)
\definecolor{lightback}{HTML}{F4E0BF}
     \definecolor{back}{HTML}{F3C591}
\definecolor{lightline}{HTML}{FCAF5F}
     \definecolor{line}{HTML}{ED7900}
% Colores para portada
\definecolor{epsc:oscuro}{HTML}{280091}
 \definecolor{epsc:medio}{HTML}{4C5CC5}
 \definecolor{epsc:claro}{HTML}{3FCFD5}
 \definecolor{epsc:verde}{HTML}{00B299}

\usepackage{tikz} % used in cover to place images

\usepackage{datetime} % allow formal date format
% "Month, YEAR" date format, in spanish with the month uppercased not interfering other dates
\newcommand\Monthname[1][EMPTY]{
  \ifnum1=#1Enero\else
  \ifnum2=#1Febrero\else
  \ifnum3=#1Marzo\else
  \ifnum4=#1Abril\else
  \ifnum5=#1Mayo\else
  \ifnum6=#1Junio\else
  \ifnum7=#1Julio\else
  \ifnum8=#1Agosto\else
  \ifnum9=#1Septiembre\else
  \ifnum10=#1Octubre\else
  \ifnum11=#1Noviembre\else
  \ifnum12=#1Diciembre\else
  \fi\fi\fi\fi\fi\fi\fi\fi\fi\fi\fi\fi
}
\newdateformat{monthyeardate}{%
  \Monthname[\THEMONTH], \THEYEAR}


% FONTs
\usepackage[T1]{fontenc}
\usepackage{textcomp}           % Needed for new symbols like € ?

\usepackage[scaled]{berasans} % Font for the cover similar to Vera 33
%\renewcommand*\familydefault{\sfdefault}  %% To use as the base font of the document is to be sans serif


% PAGE STYLE
\usepackage[twoside,bindingoffset=0mm,headheight=16pt,margin=25mm]{geometry} %,verbose,showframe
%\usepackage{fancyhdr} % Encabezados
%\pagestyle{fancy}
% XXX Evitar el binding en la portada pero no en el resto del documento

% TABLES AND FIGURES
\usepackage{graphicx}
\graphicspath{ {img/} }

% SUBRALLADOS
\usepackage[normalem]{ulem}



\begin{document}
%\frontmatter

%------------- Cover --------------
\thispagestyle{empty}

% Backgroud images
\begin{tikzpicture}[remember picture, overlay]
  % Top
  \node [anchor=north east, inner sep=0pt]  at (current page.north east)
     {\includegraphics[height=6cm]{topRightCorner.pdf}};
  % Bottom
  \node [anchor=south west, inner sep=0pt]  at (current page.south west)
     {\includegraphics[height=6cm]{bottomLeftCorner.pdf}};
  \node (uco) [anchor=south east, inner sep=0pt, xshift=-10mm, yshift=10mm]  at (current page.south east)
        {\includegraphics[height=2cm]{uco_debajo.pdf}};
  \node [anchor=south east, inner sep=0pt, xshift=-10mm]  at (uco.south west)
% Uncomment the chosen logo and comment the others:
        {\includegraphics[height=2cm]{emblema-ing-informatica.pdf}};
%        {\includegraphics[height=2cm]{emblema-ing-industrial.pdf}};
%        {\includegraphics[height=2cm]{emblema-ing-tec-industrial.pdf}};
\end{tikzpicture}


\begin{center}
\fontfamily{\sfdefault}\selectfont
\vspace*{2cm}

\vfill
\vfill
\includegraphics[width=12.5cm]{LogotipoEPSC.pdf}
\vfill
\vfill

\large\textbf{\color{epsc:medio}
TRABAJO FIN DE GRADO EN INGENIERÍA INFORMÁTICA
}
\vfill

\Huge\textbf{\color{epsc:oscuro}
  Plataforma web de gestión de comunicados
}
\vfill
\vfill

\Large\textbf{\color{epsc:verde}
  MANUAL DE USUARIO
}
\vfill
\vfill

\large{\color{epsc:oscuro}Autor}\\
\textbf{\color{epsc:medio}{ Eduardo Arroyo Ramírez }}
\vfill

\large{\color{epsc:oscuro} Director }\\
\textbf{\color{epsc:medio} Domingo Ortiz Boyer}
\vfill



\textbf{\color{epsc:verde} \monthyeardate\today}
\vfill
\vfill
\vspace{2.7cm}
\end{center}

%-------------------------------------------------------------------------------
\clearpage

\thispagestyle{empty}
\pagecolor{white}

\cleardoublepage
%-------------------------------------------------------------------------------
\frontmatter
\tableofcontents
\listoffigures
\listoftables

\cleardoublepage
\mainmatter
\chapter{Introducción}
\section{Definición del problema}
La \emph{Comunicación Corporativa} se define como <<la totalidad de los recursos de comunicación de los que dispone una organización para llegar efectivamente a sus públicos>> \cite{riel_2001}. Dentro de este concepto, se distingue entre la \emph{Comunicación Corporativa Interna} y la \emph{Comunicación Corporativa Externa} siendo la primera <<todo lo relativo a la conexión requerida entre los miembros de una determinada estructura para acometer una metas comunes>> y la segunda, <<la vinculación de la organización con el entorno en el que desarrolla sus actividades, con el fin de alcanzar un determinado nivel de rentabilidad económica y social>> \cite{castro_2007}. Desde el ámbito empresarial y administrativo se está dando mayor y mayor importancia a la comunicación ya que ésta juega un papel fundamental en la consecución de sus objetivos.

Este tipo de comunicación ha encontrado en Internet uno de sus canales principales y tanto empresas como instituciones ejercen hoy en día su presencia a través de webs corporativas, blogs, redes sociales o incluso sus propias aplicaciones corporativas para móviles, que han proliferado enormemente durante los últimos años \cite{playstore} \cite{appstore}. Todas estas herramientas permiten no sólo la publicación de comunicados sino también la programación temporal, seguimiento y medición de su impacto.

Aunque las redes sociales son uno de los medios preferidos por las empresas \cite{linkedin} y administraciones \cite{grande2015casos} como canal de comunicación debido a la gran acogida que tienen entre la población, las aplicaciones móviles son una alternativa muy a tener en cuenta ya que cubren la muy frecuente necesidad de incluir otros servicios específicos de la organización como por ejemplo:
\begin{itemize}
    \item Sistema de reservas (mesas en un restaurante, salas en un gimnasio\dots).
    \item Encuestas (toma de decisiones en una comunidad de vecinos, sondeo de la clientela de un bar para organizar un concierto en directo\dots)
    \item Sistema de geolocalización en tiempo real de un desfile (¿por dónde está pasando ahora mismo la cabalgata de reyes?).
\end{itemize}
Además, las aplicaciones salvan algunos de los inconvenientes de las redes sociales como la cesión de datos, la alta exposición a la crítica o la obligación de adaptarse al diseño y las reglas de negocio impuestas por un tercero.

\section{Antecedentes}
\subsection{Antecedentes en comunicación corporativa}

\subsection{Antecedentes en aplicaciones móviles corporativas}
% \cleardoublepage
% \chapter{Definición del problema}
% \cleardoublepage
% \chapter{Antecedentes}
% \cleardoublepage
% \chapter{Objetivos}
El objetivo general de este trabajo es la creación de una plataforma para la difusión de comunicados por parte de organizaciones independientes que facilite la medida de su impacto.

Para la consecución de dicho objetivo, se establecen los siguientes objetivos específicos:

\section{Plataforma de publicación y seguimiento de comunicados}
Es necesario dotar al sistema de una herramienta que permita a los editores elaborar, programar, publicar y consultar las métricas de impacto de los comunicados. Para ello se desarrollará una aplicación web con acceso individualizado que cubrirá las siguientes necesidades según el tipo de usuario:
\begin{itemize}
    \item Consumidores (acceso público)
    \begin{itemize}
        \item Consulta pública de comunicados de la organización.
    \end{itemize}
    \item Editores (acceso restringido por organización)
    \begin{itemize}
        \item Administración de comunicados.
        \item Consulta de métricas de impacto de comunicaciones.
        \item Administración de tipos de registro independientes como categorías de comunicaciones, teléfonos, ubicaciones, etc. de la organización.
    \end{itemize}
    \item Administradores (responsables del mantenimiento del servicio)
    \begin{itemize}
        \item Administración de organizaciones.
        \item Administración de funcionalidades disponibles para cada organización.
        \item Administración de usuarios.
        \item Configuración del servicio.
    \end{itemize}
    \item Organizaciones
    \begin{itemize}
        \item Proporcionar una API pública para la consulta de comunicados mediante una herramienta de terceros.
    \end{itemize}
\end{itemize}

\section{Aplicación móvil tipo}
El sistema contará con una herramienta que permita a los consumidores acceder a los comunicados de una organización para consultarlos desde su teléfono móvil. Para ello se elaborará una aplicación móvil tipo. Esta aplicación tipo será precursora de las diferentes aplicaciones cliente, una por cada organización registrada en el sistema.

Además de a través de la aplicación móvil, los usuarios consumidores podrán consultar los comunicados a través de la parte publica del apartado web de su organización.
% \cleardoublepage
% \chapter{Recursos}
En esta sección se clasificarán y detallarán los recursos necesarios para la realización del trabajo.

\section{Hardware}
Los recursos hardware son las herramientas físicas empleadas para el desarrollo y la puesta en marcha del proyecto.
\begin{itemize}
    \item Ordenador PC Compatible con procesador Intel i7 y 8GB de memoria RAM.
    \item Teléfono móvil Android.
    \item Conexión a Internet de banda ancha.
\end{itemize}

\section{Software}
Bajo esta denominación se engloban las herramientas software como el sistema operativo, los entornos de programación, bibliotecas, herramientas de documentación, etc.
\begin{itemize}
    \item S. O. Windows 10 Home Premium
    \item Microsoft Visual Studio Professional 2017
    \item Telerik UI para ASP.NET MVC
    \item Microsoft SQL Server Express 2016
    \item Titanium Studio
    \item Internet Information Server 7
    \item Navegadores:
        \begin{itemize}
            \item Google Chrome
            \item Mozilla Firefox
            \item Microsoft Edge
        \end{itemize}
\end{itemize}

\section{Humanos}
En este apartado se identifica a las personas involucradas en el trabajo.
\begin{itemize}
    \item Alumno: Eduardo Arroyo Ramírez
    \item Director: Domingo Ortiz Boyer
\end{itemize}
% \cleardoublepage
% \chapter{Aplicación WEB}
% \cleardoublepage
% \chapter{Servicio WEB}
% \cleardoublepage
% \chapter{Aplicación móvil}
% \cleardoublepage
% \chapter{Seguridad}
% \cleardoublepage
% \chapter{Futuras mejoras}
% \cleardoublepage
% \chapter{Conclusiones}

% BIBLIOGRAFÍA -----------------------------------------------------------------
\bibliographystyle{ieeetr}
\renewcommand{\refname}{Bibliografía}
\bibliography{00-manual-usuario}
\addcontentsline{toc}{part}{Bibliografía}
\clearpage
% ------------------------------------------------------------------------------

\end{document}